\chapter{Lịch sử mua hàng --- Giao hàng --- Báo cáo}

Chương này gồm \textbf{11 đường dẫn} gom thành \textbf{8 màn}, thuộc ba mô-đun:
\rt{wujia_portal_purchase_history} (5), \rt{wujia_portal_delivery} (4),
\rt{wujia_portal_report} (2). Ba phân hệ này đều \textbf{chỉ đọc} --- không màn nào tạo hay sửa
dữ liệu nghiệp vụ, trừ việc sinh tệp tải về.

\textbf{Ba điểm cần biết trước khi đọc:}

\begin{enumerate}[leftmargin=*]
  \item \textbf{Ba phân hệ xác định ``cửa hàng đang xem'' theo hai cách khác nhau.} Lịch sử mua
        hàng \textbf{bắt buộc} phải chọn được cửa hàng, chưa chọn thì báo lỗi và không hiện gì.
        Giao hàng và Báo cáo thì \textbf{cộng gộp mọi cửa hàng} khi chưa chọn. Đây là khác biệt
        có thật trong code, không phải mô tả nhầm --- xem mục 7 của màn Lịch sử.
  \item \textbf{Trạng thái đơn mà cửa hàng nhìn thấy là trạng thái GỘP.} Odoo chỉ có bốn trạng
        thái đơn (\emph{Nháp}, \emph{Đã gửi}, \emph{Đã xác nhận}, \emph{Đã huỷ}). Portal thêm
        hai nhãn \emph{Đang giao} và \emph{Hoàn tất} \textbf{suy ra từ chuyến giao} và cho đè
        lên nhãn \emph{Đã xác nhận}. Đây là phương án (a) chủ dự án chốt 03/08: cửa hàng chỉ
        nhìn một cột trạng thái duy nhất.
  \item \textbf{Đơn Đã huỷ bị giấu khỏi phân hệ Lịch sử.} Trên UAT hiện có \uat{16} đơn đã huỷ
        trên tổng \uat{47} --- tức \textbf{một phần ba số đơn không xuất hiện} ở màn Lịch sử.
        Nhưng chúng \textbf{vẫn có} trong tệp Excel của phân hệ Báo cáo. Xem mục 7 màn Báo cáo.
\end{enumerate}

\section{Màn Lịch sử mua hàng}
\rt{/portal/purchase-history} · \rt{/portal/purchase-history/results}
\medskip

\mvao

Mục \emph{Lịch sử mua hàng} trên thanh điều hướng. Đường dẫn \rt{/results} không phải màn
riêng --- đó là phần kết quả được gọi ngầm khi người dùng đổi bộ lọc, để không phải tải lại cả
trang. Hai đường dẫn dùng chung một bộ dữ liệu nên mọi luật dưới đây đúng cho cả hai.

Ngoài ra có \textbf{một đường dẫn cũ} \rt{/portal/purchase_history} (gạch dưới, từ bản v14) được
giữ lại để chuyển hướng vĩnh viễn sang đường dẫn mới --- liên kết cũ trong email hay dấu trang
của cửa hàng không bị chết.

\mai

\begin{itemize}[leftmargin=*]
  \item Phải đăng nhập \textbf{và phải chọn được cửa hàng}. Chưa chọn (tài khoản nhiều cửa hàng
        chưa bấm chọn) $\rightarrow$ màn hiện đúng một dòng:
        \emph{``Không xác định được cửa hàng đang thao tác. Vui lòng chọn lại cửa hàng.''} và
        \textbf{không chạy truy vấn nào}.
  \item Chỉ thấy đơn của \textbf{đúng một cửa hàng} đang chọn --- khác với Trang chủ (cộng gộp)
        và khác với Giao hàng / Báo cáo.
  \item \textbf{Vai trò không giới hạn gì}: Nhân viên xem được toàn bộ lịch sử đơn và số tiền
        của cửa hàng, ngang Chủ tiệm.
  \item Đọc bằng quyền hệ thống; giới hạn duy nhất là điều kiện \rt{franchise_id} trong bảng
        dưới.
\end{itemize}

\mhien

\begin{hienbang}
Bảng đơn hàng &
  \rt{sale.order} &
  Đơn của \textbf{đúng cửa hàng đang chọn} và \textbf{không phải đơn huỷ}
  (\rt{state != 'cancel'}). Lấy \textbf{cả} đơn đặt qua portal lẫn đơn Ngô Gia tạo hộ trong
  ERP. &
  \emph{Ngày tạo} mới nhất trước; 10 đơn/trang &
  Cầu Giấy \uat{8} · Đống Đa \uat{5} · TP HCM Q1 \uat{18} · TEST \uat{0} \\ \addlinespace
Cột \emph{Trạng thái} &
  \rt{sale.order.state} + \rt{stock.picking.batch.delivery_batch_status} &
  Năm nhãn: \emph{Chờ xác nhận} (nháp), \emph{Đã gửi}, \emph{Đã xác nhận}, và hai nhãn đè từ
  chuyến giao: \emph{Đang giao}, \emph{Hoàn tất}. Đơn đã xác nhận mà chuyến đang chạy thì hiện
  \emph{Đang giao}; chuyến đã xong thì hiện \emph{Hoàn tất}. &
  --- &
  \uat{17} Chờ xác nhận · \uat{0} Đã gửi · \uat{14} Đã xác nhận \\ \addlinespace
Cột \emph{Số dòng sản phẩm} &
  \rt{sale.order.line} &
  Đếm dòng hàng thật, \textbf{bỏ} dòng tiêu đề/ghi chú &
  --- & --- \\ \addlinespace
Cột \emph{Người đặt} &
  \rt{sale.order.portal_requester_user_id} &
  Đơn đặt qua portal $\rightarrow$ tên người bấm gửi. Đơn Ngô Gia tạo trong ERP $\rightarrow$
  nhãn cố định \emph{``Ngô Gia tạo đơn''}; \textbf{cố ý không} lộ tên nhân sự nội bộ. &
  --- &
  \uat{23} đơn portal / \uat{24} đơn tạo từ ERP \\ \addlinespace
Cột \emph{Chuyến giao} &
  \rt{stock.picking.batch} &
  Mã chuyến và trạng thái chuyến (nhãn tiếng Việt cố định trong portal, không đổi theo ngôn
  ngữ người dùng) &
  --- &
  \uat{3}/47 đơn đã được xếp chuyến \\ \addlinespace
Cột \emph{Tổng tiền} &
  \rt{sale.order.amount_total} &
  Đọc thẳng số Odoo đã tính. Ký hiệu tiền lấy theo \textbf{loại tiền của chính đơn}. &
  --- &
  \uat{47}/47 đơn đang ghi bằng \textbf{USD} \\
\end{hienbang}

\mloc

\begin{itemize}[leftmargin=*]
  \item \textbf{Từ khoá} --- chỉ tìm theo \textbf{mã đơn}, khớp một phần, không phân biệt hoa
        thường. Không tìm theo tên sản phẩm, không tìm theo người đặt.
  \item \textbf{Khoảng ngày} --- lọc theo \emph{ngày tạo đơn}, quy đổi đúng theo múi giờ tài
        khoản (``từ ngày X'' tính từ 00:00 giờ địa phương). Khoảng ngày \textbf{ngược} thì
        \textbf{không chạy truy vấn}, giữ nguyên hai ô đã nhập và báo ngay tại thanh lọc
        \emph{``Từ ngày không được lớn hơn Đến ngày''} --- cố ý không để màn rơi vào trạng thái
        ``Chưa có đơn hàng'' làm người dùng tưởng mất dữ liệu.
  \item \textbf{Nút nhanh} --- \emph{Tháng này} / \emph{Tháng trước} tự điền hai ô ngày và
        \textbf{ghi đè} ngày người dùng vừa gõ.
  \item \textbf{Trạng thái} --- một nhãn tại một thời điểm, chọn trong đúng năm nhãn kể trên.
        Giá trị lạ bị bỏ qua (coi như không lọc).
  \item \textbf{Số đơn mỗi trang} --- chỉ nhận \textbf{10, 20, 50}; số khác thì lặng lẽ dùng 10.
\end{itemize}

\mkiem{một dòng đơn}

Không kiểm gì ở màn danh sách --- sang màn chi tiết, và màn chi tiết \textbf{kiểm lại quyền từ
đầu} (xem màn sau).

\mghi

\textbf{Không ghi gì.} Màn chỉ đọc.

\mhoi

\begin{ask}
\begin{enumerate}[leftmargin=*]
  \item \textbf{Ba phân hệ hiểu ``cửa hàng đang xem'' khác nhau.} Lịch sử mua hàng bắt buộc
        chọn cửa hàng; Giao hàng và Báo cáo thì cộng gộp mọi cửa hàng khi chưa chọn; Trang chủ
        cũng cộng gộp. Người quản lý 3 cửa hàng vì vậy thấy Trang chủ báo 31 đơn nhưng màn Lịch
        sử lại báo lỗi ``chưa chọn cửa hàng''. Anh chốt giúp một cách thống nhất cho cả portal.
  \item \textbf{Đơn Đã huỷ bị giấu hoàn toàn.} Trên UAT là \uat{16}/47 đơn. Cửa hàng không có
        cách nào tra lại đơn mình đã bị huỷ, kể cả khi biết mã đơn. Anh muốn (a) giữ nguyên,
        (b) thêm một nhãn lọc \emph{Đã huỷ} để tra khi cần?
  \item \textbf{Tìm kiếm chỉ theo mã đơn.} Cửa hàng thường nhớ ``đơn có món Hồng Trà'' chứ ít
        khi nhớ mã. Có cần tìm theo tên sản phẩm trong đơn không?
  \item \textbf{Nhân viên xem được toàn bộ số tiền các đơn.} Xác nhận đúng ý anh chứ, hay muốn
        giới hạn cột \emph{Tổng tiền} cho Chủ tiệm/Quản lý?
\end{enumerate}
\end{ask}

\section{Màn Chi tiết đơn hàng}
\rt{/portal/purchase-history/<int:order_id>}
\medskip

\mvao

Bấm một dòng ở màn Lịch sử mua hàng. Trên máy tính, đây cũng là nơi người dùng được đưa tới
ngay sau khi gửi đơn thành công (xem mục 3.4).

\mai

Phải đăng nhập, đã chọn cửa hàng, và đơn phải \textbf{thuộc đúng cửa hàng đó} và
\textbf{không phải đơn huỷ}. \textbf{Quyền được kiểm lại từ đầu ở mỗi lần mở}, không tin việc
người dùng vừa bấm từ danh sách --- chống trường hợp sửa số đơn trên thanh địa chỉ để xem đơn
cửa hàng khác. Sai điều kiện thì hiện đúng một câu: \emph{``Không tìm thấy đơn hàng hoặc bạn
không có quyền xem đơn hàng này.''} --- cùng một câu cho cả hai trường hợp, \textbf{cố ý không}
cho biết đơn đó có tồn tại hay không.

\mhien

\begin{hienbang}
Khối đầu đơn &
  \rt{sale.order} &
  Mã đơn, ngày tạo, ngày đặt, trạng thái (gộp như màn danh sách), tổng tiền, mã + tên cửa hàng,
  người đặt, số dòng sản phẩm &
  --- & --- \\ \addlinespace
Khối chuyến giao &
  \rt{stock.picking.batch} &
  Chỉ hiện khi đơn \textbf{đã được xếp chuyến}: mã chuyến, trạng thái chuyến, giờ khởi hành
  (\textbf{thực tế} nếu đã đi, chưa đi thì \textbf{dự kiến}), ghi chú giao hàng &
  --- &
  \uat{3}/47 đơn có chuyến $\Rightarrow$ đa số đơn không có khối này \\ \addlinespace
Bảng dòng sản phẩm &
  \rt{sale.order.line} &
  Dòng hàng thật của đơn, \textbf{bỏ} dòng tiêu đề/ghi chú. Mỗi dòng: tên sản phẩm, quy cách
  đóng gói, đơn vị tính, số lượng, chiết khấu, \textbf{đơn giá đã gồm thuế}, thành tiền. &
  Theo thứ tự dòng trong đơn &
  --- \\ \addlinespace
Ghi chú đơn &
  \rt{sale.order.portal_note} &
  Ghi chú người dùng đã gõ ở giỏ hàng lúc gửi đơn &
  --- & --- \\
\end{hienbang}

\textbf{Về con số đơn giá:} đơn giá hiển thị được tính lại qua bộ tính thuế của Odoo cho
\textbf{đúng một đơn vị}, \textbf{không} lấy thành tiền chia cho số lượng. Phép chia đó chỉ
đúng khi thuế là phần trăm thuần --- sai với thuế cố định, sai khi một dòng chịu nhiều thuế, và
sai khi làm tròn ở loại tiền có phần lẻ. Thành tiền thì đọc thẳng số của Odoo.

\mloc

Không có bộ lọc.

\mkiem{Tải PDF}

Xem màn tiếp theo.

\mghi

Không ghi gì.

\mhoi

\begin{ask}
Màn chi tiết \textbf{không hiện lịch sử thay đổi} của đơn --- cửa hàng không thấy được đơn bị
sửa số lượng hay bị đổi giá lúc nào, do ai. Anh có cần một khối ``Nhật ký đơn'' không?
\end{ask}

\section{Tải đơn hàng dạng PDF}
\rt{/portal/purchase-history/<int:order_id>.pdf}
\medskip

\mvao

Nút \emph{Tải PDF} ở màn chi tiết đơn.

\mai

Cùng điều kiện với màn chi tiết: đúng cửa hàng đang chọn, không phải đơn huỷ. Không đạt thì trả
về lỗi \emph{không tìm thấy} chứ không chuyển hướng.

\mhien

Không phải màn --- trả thẳng về tệp PDF, tên tệp là mã đơn.

\mloc

Không có tham số nào.

\mkiem{Tải PDF}

\begin{kiembang}
1 & Đã chọn được cửa hàng & Lỗi \emph{không tìm thấy} \\
2 & Đơn thuộc đúng cửa hàng đang chọn và chưa huỷ & Lỗi \emph{không tìm thấy} \\
3 & Mẫu in đơn hàng của Odoo còn tồn tại trong hệ thống & Lỗi \emph{không tìm thấy} \\
\end{kiembang}

\mghi

Không ghi gì vào dữ liệu nghiệp vụ.

\mhoi

\begin{ask}
\textbf{Tệp PDF đang dùng mẫu in Báo giá/Đơn bán mặc định của Odoo}, chưa phải mẫu có nhận diện
Wujia Tea (logo, thông tin công ty, điều khoản). Anh có muốn đặt mẫu riêng cho cửa hàng nhượng
quyền không --- và nếu có thì PDF nên hiện những mục gì?
\end{ask}

\section{Màn Theo dõi giao hàng}
\rt{/portal/delivery} · \rt{/portal/delivery/results}
\medskip

\mvao

Mục \emph{Giao hàng} trên thanh điều hướng. Đường dẫn \rt{/results} là phần kết quả gọi ngầm khi
bấm chip lọc hoặc đổi bộ lọc.

\mai

\begin{itemize}[leftmargin=*]
  \item Phải đăng nhập. \textbf{Chưa chọn cửa hàng thì vẫn xem được} --- lấy chuyến của
        \textbf{mọi} cửa hàng tài khoản thuộc (khác hẳn màn Lịch sử mua hàng).
  \item Tài khoản không thuộc cửa hàng nào $\rightarrow$ màn trống, không báo lỗi.
  \item Một chuyến xe có thể chở hàng cho \textbf{nhiều cửa hàng}. Màn chỉ hiện \textbf{phần
        phiếu giao thuộc cửa hàng của người xem}; phiếu của cửa hàng khác trong cùng chuyến
        \textbf{không} hiện.
  \item Vai trò không giới hạn gì.
\end{itemize}

\mhien

\begin{hienbang}
Danh sách chuyến &
  \rt{stock.picking.batch} &
  Chuyến có phiếu giao thuộc cửa hàng đang xét --- nhận theo \textbf{hai đường}: phiếu gắn
  thẳng cửa hàng, \textbf{hoặc} phiếu gắn đơn hàng của cửa hàng. &
  \emph{Giờ khởi hành dự kiến} muộn nhất trước; 20 chuyến mỗi trang &
  \uat{3} chuyến toàn hệ, \uat{2} chuyến có cửa hàng nhìn thấy (một chuyến chở chung TP HCM
  Q1 + Đống Đa); Cầu Giấy \uat{0} \\ \addlinespace
Chip lọc trạng thái &
  \rt{stock.picking.batch.delivery_batch_status} &
  Bốn chip: \emph{Tất cả} / \emph{Sắp giao} (nháp, đã gán xe, đang chất hàng) / \emph{Đang giao}
  / \emph{Đã giao}. Số trên chip đếm trên \textbf{chính tập đã lọc} theo ngày và từ khoá, nên
  chọn một chip không làm các chip khác về 0. &
  --- &
  \uat{1} nháp · \uat{1} đã giao · \uat{1} huỷ chuyến \\ \addlinespace
Giờ khởi hành trên thẻ &
  \rt{stock.picking.batch} &
  Đã xuất phát thì hiện \textbf{giờ thực tế} kèm nhãn \emph{Xuất phát (thực tế)}; chưa đi thì
  hiện \textbf{giờ dự kiến} kèm nhãn \emph{Xuất phát (dự kiến)}. &
  --- & --- \\ \addlinespace
Xe và tài xế &
  \rt{stock.picking.batch.vehicle_id} &
  Tên xe, tên tài xế, biển số. Thiếu thì hiện dấu ``---''. &
  --- & --- \\ \addlinespace
Đơn liên quan trên thẻ &
  \rt{sale.order} &
  Mã các đơn thuộc cửa hàng người xem trong chuyến đó, rút gọn dạng \emph{S00035, S00036 +3} &
  --- & --- \\
\end{hienbang}

\textbf{Lưu ý về chip ``Huỷ chuyến'':} trạng thái \emph{Hủy chuyến} \textbf{không có chip
riêng}. Chuyến bị huỷ vẫn nằm trong danh sách khi chọn \emph{Tất cả}, nhưng không lọc riêng
được. Trên UAT đang có \uat{1} chuyến như vậy.

\mloc

\begin{itemize}[leftmargin=*]
  \item \textbf{Từ khoá} --- tìm theo \textbf{mã chuyến} \textbf{hoặc} \textbf{mã đơn hàng}
        trong chuyến.
  \item \textbf{Khoảng ngày} --- lọc theo \emph{giờ khởi hành dự kiến}, quy đổi theo múi giờ
        tài khoản. Ngày ngược $\rightarrow$ không chạy truy vấn, báo tại thanh lọc.
        \textbf{Lưu ý}: chuyến \textbf{chưa đặt lịch khởi hành} sẽ bị loại khỏi kết quả khi
        dùng bộ lọc ngày.
  \item \textbf{Chip trạng thái} --- một chip tại một thời điểm.
  \item \textbf{Số chuyến mỗi trang} --- chỉ nhận \textbf{10, 20, 50}; mặc định 20.
  \item Có tham số kỹ thuật \rt{_preview} dùng để đội phát triển xem thử bốn trạng thái màn
        (đang tải / lỗi / rỗng / có dữ liệu). \textbf{Không} phải chức năng nghiệp vụ, BA không
        cần lên task.
\end{itemize}

\mkiem{một chuyến}

Không kiểm gì ở danh sách --- sang màn chi tiết, và màn đó kiểm lại quyền.

\mghi

Không ghi gì.

\mhoi

\begin{ask}
\begin{enumerate}[leftmargin=*]
  \item \textbf{Cửa hàng đông đơn nhất lại không có chuyến nào.} Trên UAT, Cầu Giấy có \uat{8}
        đơn nhưng \uat{0} chuyến giao; cả hệ thống mới \uat{3}/47 đơn được xếp chuyến. Nghĩa là
        màn Giao hàng gần như trống với phần lớn cửa hàng. Đây là do dữ liệu UAT chưa đủ, hay
        nghiệp vụ thật cũng chỉ một phần đơn đi qua chuyến?
  \item \textbf{Chuyến chưa đặt lịch bị mất khi lọc theo ngày.} Người dùng lọc ``tháng này'' sẽ
        không thấy chuyến chưa có giờ khởi hành, dù chúng đang chờ giao. Anh muốn giữ nguyên
        hay luôn hiện chúng?
  \item \textbf{Không lọc riêng được chuyến bị huỷ.} Có cần thêm chip \emph{Huỷ chuyến} không?
\end{enumerate}
\end{ask}

\section{Màn Chi tiết chuyến giao}
\rt{/portal/delivery/<int:batch_id>}
\medskip

\mvao

Bấm một thẻ chuyến ở màn Theo dõi giao hàng.

\mai

\begin{itemize}[leftmargin=*]
  \item Phải đăng nhập và thuộc ít nhất một cửa hàng.
  \item Chuyến phải có \textbf{phiếu giao gắn thẳng cửa hàng} của người xem. Không đạt thì
        \textbf{đưa về màn danh sách}, không báo lỗi.
  \item \textbf{Điều kiện này HẸP HƠN màn danh sách}: danh sách nhận chuyến theo \emph{hai}
        đường (phiếu gắn cửa hàng \textbf{hoặc} phiếu gắn đơn của cửa hàng), còn chi tiết chỉ
        nhận \emph{một} đường. Về lý thuyết có thể xảy ra trường hợp chuyến hiện ở danh sách
        nhưng bấm vào lại bị đẩy về. Trên UAT hiện \textbf{chưa} xảy ra (cả hai cách đều ra
        \uat{2} chuyến), nhưng đang có \uat{1}/18 phiếu giao chưa gắn cửa hàng --- đúng loại dữ
        liệu sẽ kích hoạt lệch này. Xem mục 7.
\end{itemize}

\mhien

\begin{hienbang}
Khối đầu chuyến &
  \rt{stock.picking.batch} &
  Mã chuyến, trạng thái, giờ khởi hành (thực tế hoặc dự kiến), giờ cập nhật gần nhất, tên cửa
  hàng, kho xuất hàng &
  --- & --- \\ \addlinespace
Thanh tiến trình 3 bước &
  \rt{stock.picking.batch.delivery_batch_status} &
  \emph{Đã xuất phát} $\rightarrow$ \emph{Đang giao} $\rightarrow$ \emph{Đã giao xong}. Bước
  được tô sáng suy ra từ trạng thái chuyến và giờ khởi hành thực tế. &
  --- & --- \\ \addlinespace
Danh sách sản phẩm trong chuyến &
  \rt{stock.move} &
  Gộp theo sản phẩm, cộng số lượng --- \textbf{chỉ} phần thuộc phiếu giao của cửa hàng người
  xem &
  Theo thứ tự gặp &
  --- \\ \addlinespace
Mã đơn liên quan &
  \rt{sale.order} &
  Các đơn của cửa hàng người xem nằm trong chuyến &
  --- & --- \\
\end{hienbang}

\textbf{Về giờ ở thanh tiến trình:} hai mốc \emph{Đang giao} và \emph{Đã giao xong} hiện
\textbf{giờ cập nhật bản ghi gần nhất}, không phải giờ nghiệp vụ thật của từng bước --- hệ thống
hiện chưa lưu riêng ``giờ bắt đầu giao'' và ``giờ giao xong''. Xem mục 7.

\mloc

Không có bộ lọc.

\mkiem{Thêm vào lịch}

Xem màn tiếp theo.

\mghi

Không ghi gì.

\mhoi

\begin{ask}
\begin{enumerate}[leftmargin=*]
  \item \textbf{Điều kiện xem chi tiết hẹp hơn điều kiện hiện ở danh sách.} Nên sửa cho giống
        nhau --- em ghi nhận là việc kỹ thuật, không cần BA quyết, chỉ báo để anh biết có một
        đường dẫn sẽ được vá ở cụm F.
  \item \textbf{Chưa có giờ thật cho hai bước cuối.} Thanh tiến trình đang mượn ``giờ sửa bản
        ghi'' để hiển thị. Nếu anh cần cửa hàng thấy đúng \emph{mấy giờ xe bắt đầu giao} và
        \emph{mấy giờ giao xong} thì phải bổ sung hai trường ở phân hệ giao vận --- đây là
        thay đổi có khối lượng, cần anh xác nhận trước.
\end{enumerate}
\end{ask}

\section{Tải lịch giao hàng vào ứng dụng lịch}
\rt{/portal/delivery/<int:batch_id>.ics}
\medskip

\mvao

Nút \emph{Thêm vào lịch} ở màn chi tiết chuyến. Tệp tải về mở được bằng Google Calendar,
Outlook, Lịch của iPhone.

\mai

Cùng điều kiện với màn chi tiết chuyến. Không đạt thì trả lỗi \emph{không tìm thấy}.

\mhien

Không phải màn. Tệp lịch chứa: tiêu đề \emph{``Giao hàng + mã chuyến''}, mô tả liệt kê các
phiếu giao, địa điểm là tên các cửa hàng nhận hàng.

\mloc

Không có tham số nào.

\mkiem{Thêm vào lịch}

\begin{kiembang}
1 & Tài khoản thuộc ít nhất một cửa hàng & Lỗi \emph{không tìm thấy} \\
2 & Chuyến có phiếu giao gắn cửa hàng của người xem & Lỗi \emph{không tìm thấy} \\
\end{kiembang}

\mghi

Không ghi gì.

\mhoi

\begin{ask}
\textbf{Giờ trong tệp lịch lấy từ một trường khác với giờ hiển thị trên màn.} Màn hiện
\emph{giờ khởi hành} (thực tế hoặc dự kiến), còn tệp lịch lấy \emph{ngày giờ dự kiến của phiếu
giao} và mặc định kéo dài \textbf{2 tiếng}. Hai con số này có thể lệch nhau. Anh chốt: sự kiện
trong lịch nên bắt đầu lúc nào và kéo dài bao lâu?
\end{ask}

\section{Màn Báo cáo đặt hàng}
\rt{/portal/reports/orders}
\medskip

\mvao

Mục \emph{Báo cáo} trên thanh điều hướng. \textbf{Nhân viên không nhìn thấy mục này} và vào
thẳng đường dẫn cũng bị đẩy về Trang chủ.

\mai

\begin{itemize}[leftmargin=*]
  \item Phải đăng nhập và thuộc ít nhất một cửa hàng.
  \item \textbf{Phải là Chủ tiệm hoặc Quản lý.} Nhân viên $\rightarrow$ đẩy về Trang chủ, không
        thông báo. Trên UAT có \uat{4}/10 bản ghi thành viên là Nhân viên.
  \item \textbf{Cách xét vai trò ở đây rộng hơn phạm vi dữ liệu}: hệ thống lấy vai trò
        \textbf{cao nhất trên MỌI cửa hàng} tài khoản thuộc, rồi cho xem số liệu của
        \textbf{cửa hàng đang chọn}. Nghĩa là người chỉ là Nhân viên ở cửa hàng A nhưng là Quản
        lý ở cửa hàng B vẫn xem được báo cáo của cửa hàng A. Trên UAT có đúng \uat{1} tài khoản
        rơi vào tình huống này (Quản lý ở Đống Đa, Nhân viên ở Cầu Giấy và TP HCM Q1). Xem mục 7.
  \item Chưa chọn cửa hàng $\rightarrow$ \textbf{cộng gộp mọi cửa hàng} tài khoản thuộc.
\end{itemize}

\mhien

\begin{hienbang}
4 thẻ số liệu &
  \rt{sale.order} &
  Trong khoảng ngày đang lọc, của cửa hàng đang xét: \emph{Tổng số đơn} và \emph{Tổng doanh
  thu} (\textbf{loại} đơn huỷ), \emph{Đơn đã huỷ}, \emph{Đơn hoàn thành}. &
  --- &
  \uat{31} đơn không huỷ · \uat{16} đơn huỷ · \uat{0} đơn ``hoàn thành'' \\ \addlinespace
Biểu đồ cột theo tháng &
  \rt{sale.order} &
  Số đơn và doanh thu theo từng tháng, \textbf{loại} đơn huỷ. Khoảng luôn được nới ra
  \textbf{ít nhất 12 tháng gần nhất} để biểu đồ không trống, kể cả khi người dùng lọc một
  khoảng ngắn. &
  Theo tháng tăng dần &
  --- \\ \addlinespace
Bảng Top 10 sản phẩm &
  \rt{sale.order.line} &
  Dòng hàng thuộc đơn \textbf{đã xác nhận} của cửa hàng đang xét, trong khoảng ngày. Gộp theo
  sản phẩm, cộng số lượng và thành tiền. &
  Số lượng nhiều nhất trước; 10 dòng &
  --- \\ \addlinespace
Biểu đồ tròn + bảng phân bố trạng thái &
  \rt{sale.order} &
  Đếm đơn theo từng trạng thái kèm tỷ lệ \%. \textbf{Có tính} cả đơn huỷ. &
  --- &
  \uat{17} Nháp · \uat{0} Đã gửi · \uat{14} Đã xác nhận · \uat{0} Hoàn thành · \uat{16} Đã huỷ \\
\end{hienbang}

\mloc

Hai ô \emph{Từ ngày} / \emph{Đến ngày}, lọc theo \emph{ngày đặt hàng}. Bỏ trống thì mặc định là
\textbf{từ 01/01 năm nay đến hôm nay}. Ngày ngược $\rightarrow$ giữ nguyên chữ đã gõ, báo tại
thanh lọc, mọi số liệu và biểu đồ về rỗng.

\mkiem{Tải Excel}

Xem màn tiếp theo.

\mghi

Không ghi gì.

\mhoi

\begin{ask}
\begin{enumerate}[leftmargin=*]
  \item \textbf{Thẻ ``Đơn hoàn thành'' luôn bằng 0.} Odoo 19 chỉ có bốn trạng thái đơn
        (\emph{Nháp}, \emph{Đã gửi}, \emph{Đã xác nhận}, \emph{Đã huỷ}) --- \textbf{không có}
        trạng thái \emph{Hoàn thành}. Thẻ này và lát ``Hoàn thành'' trên biểu đồ tròn vì vậy
        không bao giờ có số. Anh muốn (a) bỏ thẻ, hay (b) định nghĩa lại ``hoàn thành'' = đơn
        đã xác nhận và chuyến giao đã xong (giống nhãn \emph{Hoàn tất} ở màn Lịch sử)?
  \item \textbf{Số tiền trên báo cáo đang dán sai ký hiệu tiền tệ.} Màn Báo cáo định dạng mọi
        số tiền theo \textbf{loại tiền của công ty}, trong khi màn Lịch sử mua hàng dùng
        \textbf{loại tiền của từng đơn}. Trên UAT, công ty đang để \textbf{VND} còn cả
        \uat{47}/47 đơn ghi bằng \textbf{USD} --- nghĩa là báo cáo đang hiện \emph{con số USD
        kèm ký hiệu ₫}. Đây đúng loại lỗi WJ-ORD-025 đã xử lý ở phân hệ Đặt hàng. Em đề xuất
        sửa theo hướng dùng loại tiền của đơn; anh xác nhận giúp.
  \item \textbf{Bộ lọc ngày của Báo cáo không quy đổi múi giờ}, khác với màn Lịch sử mua hàng
        và Giao hàng (hai màn này đã quy đổi). Hệ quả: chọn cùng một khoảng ngày, hai màn có thể
        ra số đơn khác nhau ở hai ngày biên. Đây là lỗi kỹ thuật cần vá --- em báo để anh biết,
        không cần BA quyết.
  \item \textbf{Vai trò được xét trên mọi cửa hàng, nhưng dữ liệu chỉ của cửa hàng đang chọn.}
        Người là Quản lý ở một cửa hàng đọc được báo cáo doanh thu của cửa hàng mà họ chỉ là
        Nhân viên. Anh muốn siết theo đúng cửa hàng đang xem không?
  \item \textbf{Biểu đồ luôn nới ra 12 tháng} kể cả khi người dùng lọc một tuần --- nên con số
        trên biểu đồ không khớp bốn thẻ số liệu phía trên. Anh muốn biểu đồ bám đúng khoảng
        lọc, hay giữ nguyên và ghi chú ``12 tháng gần nhất''?
\end{enumerate}
\end{ask}

\section{Tải báo cáo dạng Excel}
\rt{/portal/reports/orders/export.xlsx}
\medskip

\mvao

Nút \emph{Xuất Excel} ở màn Báo cáo.

\mai

Phải là Chủ tiệm hoặc Quản lý --- nhưng \textbf{xét theo cửa hàng đang chọn}, không phải theo
mọi cửa hàng như màn Báo cáo. Nghĩa là có trường hợp \textbf{xem được màn nhưng bấm xuất Excel
thì bị đẩy về Trang chủ} mà không có lời giải thích nào. Xem mục 7.

\mhien

Không phải màn --- trả về tệp Excel một trang tính tên \emph{Orders}, tên tệp kèm khoảng ngày.
Bảy cột: \emph{Mã đơn}, \emph{Ngày đặt}, \emph{Cửa hàng}, \emph{Khách hàng}, \emph{Trạng thái},
\emph{Số dòng}, \emph{Tổng tiền}.

\mloc

Nhận cùng hai ô ngày với màn Báo cáo. Khác ở chỗ: khoảng ngày ngược \textbf{không báo lỗi} mà
bị lặng lẽ kéo về đúng một ngày.

\mkiem{Xuất Excel}

\begin{kiembang}
1 & Tài khoản thuộc ít nhất một cửa hàng & Đưa về Trang chủ \\
2 & Vai trò tại \textbf{cửa hàng đang chọn} là Chủ tiệm hoặc Quản lý &
    Đưa về Trang chủ, không thông báo \\
\end{kiembang}

\mghi

Không ghi gì.

\mhoi

\begin{ask}
\begin{enumerate}[leftmargin=*]
  \item \textbf{Tệp Excel chứa cả đơn đã huỷ, màn hình thì không.} Người dùng đối chiếu hai bên
        sẽ thấy lệch \uat{16} đơn trên dữ liệu UAT. Anh chốt: tệp Excel nên có hay không có đơn
        huỷ?
  \item \textbf{Cột \emph{Số dòng} trong Excel đếm cả dòng tiêu đề/ghi chú}, còn màn hình chỉ
        đếm dòng sản phẩm thật. Hai nơi sẽ ra số khác nhau với đơn có dòng ghi chú.
  \item \textbf{Cột \emph{Tổng tiền} trong Excel là số trần, không kèm loại tiền} --- với dữ
        liệu USD/VND lẫn lộn như hiện nay thì người đọc không biết đơn vị. Nên thêm một cột
        \emph{Loại tiền}, anh đồng ý chứ?
  \item \textbf{Điều kiện vai trò của nút Xuất khác của màn} (đã nêu ở mục 2). Cần thống nhất
        theo cách anh chốt ở câu 4 mục 5.8.
\end{enumerate}
\end{ask}

\section{Đường dẫn cũ của Lịch sử mua hàng}
\rt{/portal/purchase_history}
\medskip

Đường dẫn thời bản v14 (viết gạch dưới), được giữ lại làm chuyển hướng vĩnh viễn sang
\rt{/portal/purchase-history}. Không hiện màn, không đọc dữ liệu. BA không cần lên task.
